\DSource{0139}{50}
\hypertarget{AB01-PDF0139-H01}{}
\DLine{AB01-PDF0139-body-L001}{}{p. 76.}{\CenterLine{CAPUT XXX.}}
\vspace{13.8pt}
\DLine{AB01-PDF0139-body-L002}{5}{}{\CenterLine{De sphaeris Lunae, de inaequalitate eius motus, de accretione et demi-}}
\DLine{AB01-PDF0139-body-L003}{}{}{\CenterLine{nutione eius luminis, de causis eclipsium, atque de distantiis geo-}}
\DLine{AB01-PDF0139-body-L004}{}{}{\CenterLine{centricis, diametris et magnitudinibus amborum luminarium cum}}
\DLine{AB01-PDF0139-body-L005}{}{}{\CenterLine{Terra comparatorum.}}
\vspace{13.8pt}
\hypertarget{AB01-PDF0139-P01}{}
\DLine{AB01-PDF0139-body-L006}{10}{}{\Indent Dicit [auctor]: In motibus Lunae duae inaequalitates repertae sunt. Altera, per se simplex,}
\DLine{AB01-PDF0139-body-L007}{}{}{temporibus syzygiarum, quae ob medium Solis et Lunae iter iuxta locum Lunae in epicyclo}
\DLine{AB01-PDF0139-body-L008}{}{}{fiunt, apparet. Altera per elongationes Lunae a Sole deprehenditur, et primae inaequalitati ita}
\DLine{AB01-PDF0139-body-L009}{}{}{adiungitur ut simul unam rem efficiant, sicut demonstrationibus geometricis cognoscitur.}
\hypertarget{AB01-PDF0139-P02}{}
\DLine{AB01-PDF0139-body-L010}{}{}{\Indent Luna quatuor sphaeras [seu circulos] habere cogitetur; quorum primus concentricus et}
\DLine{AB01-PDF0139-body-L011}{15}{}{similis sit eclipticae, eodemque motu sub ea absque ulla deviatione moveatur (1) --- Secundus ab}
\DLine{AB01-PDF0139-body-L012}{}{}{eo in septentriones et meridiem versus deflectat, sed eandem magnitudinem idemque centrum}
\DLine{AB01-PDF0139-body-L013}{}{}{ac ille habeat; maxima eius declinatio in utramque plagam versus sit fere $5^\circ$, quae est maxima}
\DLine{AB01-PDF0139-body-L014}{}{}{distantia latitudinalis Lunae ab ecliptica. Motus huius circuli deflectentis [i. e. obliqui] quotidie}
\DLine{AB01-PDF0139-body-L015}{}{}{sit 3 minutorum fere in partem contrariam successioni signorum; et hic est motus duorum}
\DLine{AB01-PDF0139-body-L016}{20}{}{nodorum, quorum alter, a quo Luna in latitudine ad septentriones transit, {\itshape Caput} vocatur, alter}
\DLine{AB01-PDF0139-body-L017}{}{}{vero, a quo Luna ad meridiem transit, {\itshape Cauda}. Hi nodi sunt loca ubi circulus deflectens cir-}
\DLine{AB01-PDF0139-body-L018}{}{}{culum pareclipticum secat. --- Intra circulum deflectentem [seu obliquum] tertius reperiatur,}
\DLine{AB01-PDF0139-body-L019}{}{}{qui excentricus sit, a deflectente pendat, eumque in puncto omnium altissimo, {\itshape apogeo} appel-}
\DLine{AB01-PDF0139-body-L020}{}{}{lato, contingat; ipse quotidie intra circulum deflectentem versus partem contrariam successioni}
\DLine{AB01-PDF0139-body-L021}{25}{}{signorum $11^\circ12'$ circiter (2) moveatur. --- Quartus circulus sit epicyclus Lunae, cuius cen-}
\DLine{AB01-PDF0139-body-L022}{}{}{trum in circumferentia excentrici quotidie $24^\circ23'$ fere (3) ad successionem signorum mo-}
\DLine{AB01-PDF0139-body-L023}{}{p. 77.}{veatur, ut sit initium motus a puncto apogei excentrici dato cum loco medio Solis. Centrum}
\DLine{AB01-PDF0139-body-L024}{}{}{igitur epicycli bis in mense lunari in punctum apogei incidit; altera vice tempore mediae}
\DLine{AB01-PDF0139-body-L025}{}{}{coniunctionis, altera tempore oppositionis; Luna autem quotidie in circumferentia epicycli}
\DLine{AB01-PDF0139-body-L026}{30}{}{fere $13^\circ4'$ (4) versus partem contrariam successioni signorum movetur, incipiens a puncto}
\DLine{AB01-PDF0139-body-L027}{}{}{apogei quod respectu centri excentrici consideratur. Et cum centrum epicycli super circum-}
\DLine{AB01-PDF0139-body-L028}{}{}{ferentiam circuli deflectentis duobus memoratis temporibus incidat, nihil obstat quin epicycli}
\DLine{AB01-PDF0139-body-L029}{}{}{centrum quotidie super circumferentiam deflectentis $13^\circ14'$ fere moveatur; et hic est motus}
\DLine{AB01-PDF0139-body-L030}{}{}{Lunae in latitudine. Sed nodus eum versus partem contrariam successioni signorum minuit}
\DLine{AB01-PDF0139-body-L031}{35}{}{tribus minutis memoratis, quae sunt motus circuli deflectentis; qua re motus Lunae in lon-}
\DLine{AB01-PDF0139-body-L032}{}{}{gitudine ad successionem signorum remanet $13^\circ11'$ fere in die. Motus in epicyclo est primus}
\DLine{AB01-PDF0139-body-L033}{}{}{quem memoravimus.}
\hypertarget{AB01-PDF0139-P03}{}
\DLine{AB01-PDF0139-body-L034}{}{}{\Indent Ex iis quae diximus patet motui Lunae nullam variationem ob excentricum his duobus tem-}
\DLine{AB01-PDF0139-body-L035}{}{}{poribus contingere, nam in iis Luna a loco Solis medio vel huic opposito non recedit, et inae-}
\DLine{AB01-PDF0139-body-L036}{40}{}{qualitas simplex cum altera non commiscetur eo usque quo Luna a Sole non recesserit; et}
\par\vspace{6pt}\hrule height.25pt\vspace{5pt}
\noindent\begin{minipage}[t]{.48\textwidth}\vspace{0pt}\fontsize{9}{11.5}\selectfont\setlength{\GutterOffset}{0pt}
\hypertarget{AB01-PDF0139-N01}{}
\DLine{AB01-PDF0139-notes_left-L001}{}{}{\Indent (1) Eum infra {\itshape pareclipticum} voco; cfr. p. 45,}
\DLine{AB01-PDF0139-notes_left-L002}{}{}{adn. 3.}
\hypertarget{AB01-PDF0139-N02}{}
\DLine{AB01-PDF0139-notes_left-L003}{}{}{\Indent (2) Nam motus apogei = duplus motus syno-}
\DLine{AB01-PDF0139-notes_left-L004}{}{}{dicus $-$ motus in latit. $+$ motus nodorum $=24^\circ$}
\DLine{AB01-PDF0139-notes_left-L005}{}{}{$23'-13^\circ14'+0^\circ3'=11^\circ12'$.}
\end{minipage}
\hfill\begin{minipage}[t]{.48\textwidth}\vspace{0pt}\fontsize{9}{11.5}\selectfont\setlength{\GutterOffset}{0pt}
\hypertarget{AB01-PDF0139-N03}{}
\DLine{AB01-PDF0139-notes_right-L001}{}{}{\Indent (3) Est duplus motus synodicus (vel relativus}
\DLine{AB01-PDF0139-notes_right-L002}{}{}{sive dupla elongatio.}
\hypertarget{AB01-PDF0139-N04}{}
\DLine{AB01-PDF0139-notes_right-L003}{}{}{\Indent (4) Est motus diurnus anomaliae, quem \Name{Pto-}}
\DLine{AB01-PDF0139-notes_right-L004}{}{}{\Name{lemaeus} IV, 3 (ed. Halma t. I, p. 223) invenit esse}
\DLine{AB01-PDF0139-notes_right-L005}{}{}{$13^\circ3'53''56'''29^{\mathrm{IV}}38^{\mathrm{V}}38^{\mathrm{VI}}$.}
\end{minipage}
